\subsection{Distribuzione dell'impegno orario}
La seguente tabella illustra la distribuzione equa delle 93 ore individuali per ciascun componente del gruppo Git Happens tra i 6 ruoli definiti.

\vspace{2mm}
\noindent
\begin{tabularx}{\textwidth}{|X c c c c c c | c|}
\hline
\textbf{Componente} & \textbf{Resp} & \textbf{Amm} & \textbf{Analist} & \textbf{Progett} & \textbf{Programm} & \textbf{Verif} & \textbf{Totale} \\
\hline
Filippo Alessandrin & 7 & 7 & 10 & 20 & 24 & 25 & \textbf{93} \\
Cristiano Antoniazzo & 7 & 7 & 10 & 20 & 24 & 25 & \textbf{93} \\
Tommaso Caon & 7 & 7 & 10 & 20 & 24 & 25 & \textbf{93} \\
Gabriele Gambarro & 7 & 7 & 10 & 20 & 24 & 25 & \textbf{93} \\
Gianmaria Grinovero & 7 & 7 & 10 & 20 & 24 & 25 & \textbf{93} \\
Diana Grosselle & 7 & 7 & 10 & 20 & 24 & 25 & \textbf{93} \\
Antonio Hai Xu Tang & 7 & 7 & 10 & 20 & 24 & 25 & \textbf{93} \\
\hline
\textbf{TOTALE ORE} & \textbf{49} & \textbf{49} & \textbf{70} & \textbf{140} & \textbf{168} & \textbf{175} & \textbf{651} \\
\hline
\end{tabularx}
\\ \\
La ripartizione delle ore riflette il peso atteso delle attività nel contesto del progetto Sonar. A ciascun componente del gruppo sono assegnate 93 ore, per un totale di 651 ore, distribuite tra i ruoli secondo la seguente logica.
\\ \\
\textbf{Responsabile e Amministratore (49 ore ciascuno).} Si tratta di ruoli di coordinamento e supporto, che non producono direttamente il prodotto ma ne garantiscono la corretta conduzione. Comprendono la pianificazione e il coordinamento del gruppo, la gestione degli incontri quindicinali con l'azienda proponente (comprensivi della revisione dei consumi delle piattaforme Vercel e Supabase), la gestione delle comunicazioni asincrone e la redazione e manutenzione delle norme di progetto. Trattandosi di attività continue ma di impegno contenuto, il peso complessivo è limitato.
\\ \\
\textbf{Analista (70 ore).} L'attività di analisi è concentrata nelle fasi iniziali e comprende l'analisi dello stato dell'arte, l'analisi dei requisiti e la redazione dei casi d'uso in UML. Poiché il capitolato definisce già in modo articolato i requisiti obbligatori, desiderabili e opzionali, l'impegno è rivolto soprattutto alla loro interpretazione, al tracciamento e all'eventuale negoziazione con l'azienda proponente, ad esempio per la riduzione concordata dell'ambito prevista dal capitolato.
\\ \\
\textbf{Progettista (140 ore).} Il ruolo assume un peso doppio rispetto all'analisi, poiché la qualità dell'architettura è il criterio su cui Halue valuterà più severamente il progetto, in particolare la separazione netta tra acquisizione, normalizzazione, analisi e presentazione. L'attività comprende la specifica architetturale, la modellazione del grafo delle entità e delle dipendenze, la strategia di campionamento, la gestione del rate limiting, la struttura del motore di regole e la redazione degli Architecture Decision Record.
\\ \\
\textbf{Programmatore (168 ore).} Rappresenta la quota dedicata alla realizzazione del Proof of Concept e del Minimum Viable Product: scansione durevole tramite Vercel Workflows, analisi statica del tema Liquid, motore di regole estensibile, persistenza su Supabase con Row Level Security, interfaccia web e funzionalità di diff ed esportazione. Il peso, pur elevato, è contenuto dall'uso di strumenti di assistenza alla scrittura del codice, che sposta la difficoltà dalla scrittura alle decisioni progettuali.
\\ \\
\textbf{Verificatore (175 ore).} È il ruolo con il maggiore impegno orario, in coerenza con i criteri di accettazione del capitolato: copertura dei test pari ad almeno il 70\% verificata in CI, test book completo, percentuale di casi superati non inferiore al 95\%, tasso di falsi positivi delle segnalazioni ad alta severità non superiore al 10\% e test automatici sull'isolamento tra organizzazioni. L'attività comprende inoltre la verifica della documentazione e della qualità del codice.
\\ \\
La distribuzione è identica per tutti i componenti, così da garantire un carico equo e una conoscenza trasversale di tutte le parti del sistema, in linea con il requisito dell'azienda proponente secondo cui ogni parte del prodotto deve poter essere spiegata da almeno un membro del gruppo.

% ------------------------------------------------------------------------
% 4. PREVENTIVO DEI COSTI
% ------------------------------------------------------------------------
\subsection {Preventivo dei costi}
Nella seguente tabella vengono riportati i costi orari per ruolo, le ore totali previste ed il costo complessivo stimato:

\vspace{2mm}
\noindent
\begin{tabularx}{\textwidth}{|X c c c r|}
\hline
\textbf{Ruolo} & \textbf{Costo orario (\texteuro{})} & \textbf{Ore per ruolo} & \textbf{Ore per membro} & \textbf{Costo totale (\texteuro{})} \\
\hline
Responsabile & 30 & 49 & 7 & 1.470,00 \\
Amministratore & 20 & 49 & 7 & 980,00 \\
Analista & 25 & 70 & 10 & 1.750,00 \\
Progettista & 25 & 140 & 20 & 3.500,00 \\
Programmatore & 15 & 168 & 24 & 2.520,00 \\
Verificatore & 15 & 175 & 25 & 2.625,00 \\
\hline
\textbf{TOTALE} & \textbf{-} & \textbf{651} & \textbf{93} & \textbf{12.865,00} \\
\hline
\end{tabularx}

\vspace{2mm}
Il preventivo di spesa totale stimato per l'esecuzione dell'intero progetto è di \textbf{\texteuro{} 12.865,00}. La quantificazione economica è stata calcolata in conformità alle tariffe stabilite dal Regolamento di Progetto. \\
Il gruppo considera questo importo come limite massimo per l'intero progetto.